% Plantilla de informe técnico / memoria en español. Compilar con:
%   latexmk -pdf -interaction=nonstopmode -halt-on-error main.tex
\documentclass[11pt,a4paper]{report}
\usepackage[T1]{fontenc}
\usepackage[spanish,es-tabla]{babel}
\usepackage{lmodern,microtype}
\usepackage[margin=2.5cm]{geometry}
\usepackage{amsmath,amssymb}
\usepackage{graphicx,booktabs,array}
\usepackage{xcolor}
\usepackage{siunitx}
\sisetup{output-decimal-marker={,}}
\usepackage{fancyhdr}
\pagestyle{fancy}
\fancyhf{}
\fancyhead[L]{\nouppercase{\leftmark}}
\fancyfoot[C]{\thepage}
\setlength{\headheight}{14pt}
\usepackage[style=apa,backend=biber]{biblatex}
\addbibresource{referencias.bib}
\usepackage[hidelinks]{hyperref}
\usepackage[spanish,noabbrev]{cleveref}

\title{Título del informe\\[0.5em]\large Subtítulo o proyecto}
\author{Nombre Apellido\\ Departamento, Universidad}
\date{\today}

\begin{document}
\maketitle
\thispagestyle{empty}

\begin{abstract}
Resumen de 150--250 palabras: contexto, objetivo, método, resultados principales (con cifras) y conclusión.
\end{abstract}

\tableofcontents

\chapter{Introducción}
\label{cap:intro}
Contexto y motivación. Estado de la cuestión \parencite{ejemplo2023}. Objetivo del trabajo.

\chapter{Materiales y métodos}
\label{cap:metodos}
Descripción reproducible de equipos, materiales, procedimiento y análisis de datos.

\chapter{Resultados}
\label{cap:resultados}
Los resultados se resumen en la \cref{tab:resumen} y la \cref{fig:ejemplo}.

\begin{table}[htbp]
  \centering
  \caption{Resumen de resultados (media $\pm$ DE).}
  \label{tab:resumen}
  \begin{tabular}{lcc}
    \toprule
    Ensayo & Valor (\si{\mega\pascal}) & $n$ \\
    \midrule
    A & $\num{250.3} \pm \num{4.1}$ & 5 \\
    B & $\num{198.7} \pm \num{6.3}$ & 5 \\
    \bottomrule
  \end{tabular}
\end{table}

\begin{figure}[htbp]
  \centering
  % \includegraphics[width=0.75\linewidth]{figuras/ejemplo.pdf}
  \fbox{\rule{0pt}{4cm}\rule{0.6\linewidth}{0pt}}
  \caption{Pie de figura descriptivo (qué se muestra y en qué condiciones).}
  \label{fig:ejemplo}
\end{figure}

\chapter{Discusión}
Interpretación, comparación con trabajos previos y limitaciones.

\chapter{Conclusiones}
Conclusiones numeradas que respondan a los objetivos y trabajo futuro.

\printbibliography[heading=bibintoc,title={Referencias}]

\appendix
\chapter{Datos adicionales}
Material complementario.

\end{document}
